% Part: computability
% Chapter: recursive-functions
% Section: general-recursive-functions

\documentclass[../../../include/open-logic-section]{subfiles}

\begin{document}

\olfileid{cmp}{rec}{gen}
\olsection{عمومی بازگشتی فنکشن}

کُل فنکشناں دا اک سیٹ حاصل کرن دا ہور وی
طریقہ اے۔ اک کُل فنکشن $f(x,\vec z)$ نوں
\emph{باقاعدہ (regular)} کہیئے جے قدرتی عدداں
دے ہر سلسلے $\vec z$ لئی اک $x$ موجود
ہووے جیہدے لئی $f(x,\vec z) = 0$۔ دوجے
لفظاں وچ، باقاعدہ فنکشن عین اوہ فنکشن نیں
جیہناں اُتے غیر محدود تلاش لا کے آخرکار کُل
فنکشن حاصل ہووے۔ احتیاط نال غیر محدود تلاش
نوں صرف باقاعدہ فنکشناں تک محدود رکھ سکدے آں:

\begin{defn}
\ollabel{defn:general-recursive}
\emph{عمومی بازگشتی فنکشناں (general recursive
functions)} دا سیٹ قدرتی عدداں توں قدرتی
عدداں ول مختلف دلیل-گنتیاں والے فنکشناں دا
سب توں چھوٹا سیٹ اے، جیہدے وچ صفر، جانشین
تے پروجیکشن فنکشن شامل نیں، تے جیہڑا ترکیب،
بنیادی بازگشت تے \emph{باقاعدہ} فنکشناں اُتے
غیر محدود تلاش لاون تحت بند اے۔
\end{defn}

صاف اے کہ ہر عمومی بازگشتی فنکشن کُل اے۔
\olref{defn:general-recursive} تے
\olref[par]{defn:recursive-fn} دا فرق ایہہ اے
کہ دوجی تعریف وچ راہ وچ جزوی بازگشتی فنکشن
ورتن دی اجازت اے؛ شرط صرف ایہہ اے کہ آخر
وچ حاصل ہون والا فنکشن کُل ہووے۔ ایس لئی لفظ
«عمومی»، جیہڑا تاریخ توں بچیا ہویا ناں اے،
معنی نال نہیں ملدا: ظاہری طور اُتے
\olref{defn:general-recursive}،
\olref[par]{defn:recursive-fn} نالوں \emph{گھٹ}
عمومی اے۔ پر خوش قسمتی نال ایہہ فرق صرف
ظاہری اے؛ تعریفاں وکھ نیں، پر عمومی بازگشتی
فنکشناں دا سیٹ تے بازگشتی فنکشناں دا سیٹ
اکّو نیں۔

\end{document}
